% ============================================================
%  sec-4-consolidated-architecture.tex  —  3.4 Consolidated Architecture
% ============================================================

\section{Consolidated Nemotron 3 Nano Architecture}
\label{sec:consolidated-architecture}

% ── 3.4.1 ────────────────────────────────────────────────────
\subsection{End-to-End Architecture Description}
\label{subsec:e2e-architecture}

The Nemotron 3 Nano training and deployment system constitutes a
vertically integrated NVIDIA-stack pipeline spanning data ingestion,
distributed training, alignment, evaluation, and production inference.
The architecture is characterised by three structural properties.

\textbf{Stage isolation with artifact-mediated handoff.} Each stage
is an independently deployable unit communicating exclusively through
W\&B artifacts.
\textbf{Framework heterogeneity across stages.} The pipeline uses
Megatron-Bridge, NeMo-RL, and cosmos\_xenna to address different
training semantics.
\textbf{Configuration-driven architecture selection.} The recipe system
enables swapping model architectures via a single configuration field
with no code changes.

% ── 3.4.2 ────────────────────────────────────────────────────
\subsection{Training Infrastructure Summary}
\label{subsec:training-infra-summary}

\begin{table}[htbp]
  \centering
  \caption{Training infrastructure summary by stage}
  \label{tab:training-summary}
  \begin{tabularx}{\textwidth}{>{\hsize=0.6\hsize}L >{\hsize=1.2\hsize}L l l l >{\hsize=1.4\hsize}L}
    \toprule
    \tableheader{Stage & Framework & Launch & Nodes & GPUs & Container} \\
    \midrule
    Data Prep
      & cosmos\_xenna + Ray
      & ray
      & 1
      & 0
      & \path{anyscale/ray:2.49.2-py312} \\
    \rowcolor{TableRowOdd}
    Pretrain
      & Megatron-Bridge
      & torchrun
      & 2
      & 8/node
      & \path{nemo:25.11.nemotron_3_nano} \\
    SFT
      & Megatron-Bridge
      & torchrun
      & 1
      & 8/node
      & \path{nemo:26.02.super.rc1} \\
    \rowcolor{TableRowOdd}
    RL (GRPO)
      & NeMo-RL + Ray
      & ray
      & 1
      & 8/node
      & \path{nemo-rl:v0.4.0.nemotron_3_nano} \\
    \bottomrule
  \end{tabularx}
\end{table}

% ── 3.4.3 ────────────────────────────────────────────────────
\subsection{Inference Infrastructure and Dependency Summary}
\label{subsec:inference-summary}

\begin{table}[htbp]
  \centering
  \caption{Inference infrastructure options}
  \label{tab:inference-summary}
  \begin{tabularx}{\textwidth}{>{\hsize=0.6\hsize}L >{\hsize=1.4\hsize}L >{\hsize=1.0\hsize}L}
    \toprule
    \tableheader{Engine & Integration Point & Deployment Context} \\
    \midrule
    vLLM
      & RL generation backend (\texttt{policy\_generation} actor); standalone serving
      & Research and production serving \\
    \rowcolor{TableRowOdd}
    TensorRT-LLM
      & Post-training compilation of final checkpoint
      & Optimised production inference \\
    NVIDIA NIM
      & Containerised packaging of TRT-LLM engine
      & Enterprise deployment; OpenAI-compatible API \\
    \rowcolor{TableRowOdd}
    NeMo Evaluator
      & Benchmark inference during Stage 3 evaluation
      & Accuracy measurement before deployment \\
    \bottomrule
  \end{tabularx}
\end{table}


% ── 3.4.5 ────────────────────────────────────────────────────
\subsection{Concluding Remarks}
\label{subsec:ch3-conclusion}

The Nemotron 3 Nano ecosystem implements a complete, production-grade
LLM lifecycle spanning pretraining, supervised fine-tuning, reinforcement
learning alignment, evaluation, and deployment. Each stage is implemented
as a discrete, independently executable unit communicating through typed
W\&B artifacts. The training stack is exclusively NVIDIA-native: every
GPU-accelerated stage depends on \acrshort{cuda}, cuDNN, NCCL, or
TensorRT-LLM with no hardware-agnostic abstraction layer between the
framework code and NVIDIA's accelerator stack. The data preparation
pipeline, operating on CPU-only Ray nodes, is the sole stage with no
direct \acrshort{cuda} dependency.

This architecture constitutes the reference NVIDIA-based ecosystem that
will be analysed and used as the baseline for the migration study
presented in the next chapter.
